\documentclass[a4paper, serif, 10pt]{moderncv}

%% moderncv
% style and color
\moderncvstyle{banking}
\moderncvcolor{black}

% renew moderncv command to adapt for CJK colon
\renewcommand*{\cvitem}[3][.25em]{%
  \ifstrempty{#2}{}{\hintstyle{#2}：}{#3}%
  \par\addvspace{#1}}

\renewcommand*{\cvitemwithcomment}[4][.25em]{%
  \savebox{\cvitemwithcommentmainbox}{\ifstrempty{#2}{}{\hintstyle{#2}：}#3}%
  \setlength{\cvitemwithcommentmainlength}{\widthof{\usebox{\cvitemwithcommentmainbox}}}%
  \setlength{\cvitemwithcommentcommentlength}{\maincolumnwidth-\separatorcolumnwidth-\cvitemwithcommentmainlength}%
  \begin{minipage}[t]{\cvitemwithcommentmainlength}\usebox{\cvitemwithcommentmainbox}\end{minipage}%
  \hfill% fill of \separatorcolumnwidth
  \begin{minipage}[t]{\cvitemwithcommentcommentlength}\raggedleft\small\itshape#4\end{minipage}%
  \par\addvspace{#1}}

%% page layout/margins
\usepackage[top=2.5cm, bottom=2.5cm, left=1.5cm, right=1.5cm]{geometry}

\nopagenumbers{}

%% Babel config for English language
\usepackage[english]{babel}

%% fontspec
\usepackage{fontspec}

\IfFontExistsTF{Linux Libertine}{
  \setmainfont[Ligatures={TeX, Common}, Numbers=Lining, ItalicFont=Linux Libertine]{Linux Libertine}
}{}
\IfFontExistsTF{Linux Libertine O}{
  \setmainfont[Ligatures={TeX, Common}, Numbers=Lining, ItalicFont=Linux Libertine O]{Linux Libertine O}
}{}

%% CTeX
% CJK support, used to show CJK characters in the resume
%
% - fontset=none: disable builtin fontset but instead set the CJK font manually
% - heading=false: disable ctex heading
% - punct=kaiming: use kaiming punctuations styles for CJK
% - scheme=plain: use plain scheme, do not override \normalsize font size
% - space=auto: space settings for CJK characters
%
% ref:
% - http://ctan.mirrorcatalogs.com/language/chinese/ctex/ctex.pdf
\usepackage[UTF8, heading=false, punct=kaiming, scheme=plain, space=auto]{ctex}

\IfFontExistsTF{Noto Serif CJK SC}{
  \setCJKmainfont{Noto Serif CJK SC}
}{}
\IfFontExistsTF{Noto Sans CJK SC}{
  \setCJKsansfont{Noto Sans CJK SC}
}{}

%% line spacing
\usepackage{setspace}
\setstretch{1.125}

%% URL styling - use normal text instead of monospace
\urlstyle{same}

% Auto-underline all links
% Defer until \begin{document} so hyperref is already loaded
\AtBeginDocument{%
  \let\oldhref\href
  \renewcommand{\href}[2]{\oldhref{#1}{\underline{#2}}}%
}

%% Patch \httplink and \httpslink to handle full URLs
% The original moderncv macros always prepend http:// or https:// to URLs.
% This causes issues when the URL already has a protocol (e.g., https://example.com)
% resulting in malformed URLs like https://https://example.com.
% This patch checks if the URL already contains "://" and uses it directly if so.
% Note: We use \str_set:Nx (x-expansion) to expand macros like \@homepage first.
\ExplSyntaxOn
\str_new:N \g_yamlresume_url_str

\RenewDocumentCommand{\httpslink}{O{}m}{
  \str_set:Nx \g_yamlresume_url_str {#2}
  \str_if_in:NnTF \g_yamlresume_url_str {://}
    { \url{#2} }
    { \url{https://#2} }
}

\RenewDocumentCommand{\httplink}{O{}m}{
  \str_set:Nx \g_yamlresume_url_str {#2}
  \str_if_in:NnTF \g_yamlresume_url_str {://}
    { \url{#2} }
    { \url{http://#2} }
}
\ExplSyntaxOff

\name{陳嘉敏}{}
\title{資深項目經理 | 數碼轉型與產品交付}
\phone[mobile]{+852 9123 4567}
\email{chan.ka.man.pm@gmail.com}
\homepage{https://www.chan-ka-man.com}

\address{中國香港，九龍，香港，香港九龍旺角彌敦道700號}{}{}

\extrainfo{{\small \faLinkedin}\ \href{https://www.linkedin.com/in/chan-ka-man-pm}{@chan-ka-man-pm} {} {} {} • {} {} {} 
{\small \faGithub}\ \href{https://github.com/chanpmhk}{@chanpmhk}}

\begin{document}

\maketitle

\section{簡介}

\cvline{}{資深項目經理，擁有超過 10 年資訊科技及數碼轉型項目管理經驗。擅長領導跨部門團隊，從需求分析、資源規劃到交付上線，確保項目按時、按預算及符合質量標準完成。
%
持有 PMP 及 PMI-ACP 專業認證，精通敏捷與傳統項目管理方法。曾成功交付多個超過港幣 2,000 萬的大型項目，涵蓋金融科技、電訊及客戶服務平台。
%
\begin{itemize}
\item 卓越的持份者管理及溝通能力
\item 數據驅動的決策與風險管理
\item 擅長推動變更管理及團隊賦能
\end{itemize}}
\section{教育背景}

\cventry{2010年9月–2014年6月}
        {學士，資訊系統與管理，成績：3.6}
        {香港科技大學}
        {\url{https://www.ust.hk}}
        {}
        {\begin{itemize}
\item 主修資訊系統與管理，副修商業分析
\item 曾獲院長嘉許名單
\item 畢業專題為智慧城市交通管理系統
\end{itemize}
\textbf{課程}：資料庫管理系統、系統分析與設計、專案管理、企業資源規劃、商業智能、軟件工程}
\section{工作經歷}

\cventry{2021年4月至今}
        {高級項目經理}
        {香港電訊數碼有限公司}
        {\url{https://www.hkt.com}}
        {}
        {\begin{itemize}
\item 領導跨部門團隊交付多個大型數碼轉型項目，總預算超過港幣 2,000 萬
\item 制定項目路線圖、資源分配及風險管理計劃，確保按時按預算交付
\item 與產品、工程及業務團隊協作，推動敏捷開發流程及持續改進
\item 建立項目績效指標及匯報機制，定期向管理層匯報進度
\item 指導初級項目經理及項目協調員，提升團隊整體交付能力
\end{itemize}
\textbf{關鍵字}：數碼轉型、敏捷開發、風險管理、持份者管理}

\cventry{2018年7月–2021年3月}
        {項目經理}
        {新鴻基金融集團}
        {\url{https://www.shkf.com}}
        {}
        {\begin{itemize}
\item 管理金融科技項目，包括網上銀行平台及流動支付系統
\item 協調內部開發團隊與外部供應商，確保系統整合順利
\item 主持需求收集工作坊，將業務需求轉化為技術規格
\item 實施項目治理框架，降低項目延誤及超支風險
\end{itemize}
\textbf{關鍵字}：金融科技、網上銀行、供應商管理、需求分析}

\cventry{2014年8月–2018年6月}
        {助理項目經理}
        {電訊盈科有限公司}
        {\url{https://www.pccw.com}}
        {}
        {\begin{itemize}
\item 協助管理客戶關係管理系統升級項目，涵蓋 10 個業務部門
\item 負責項目文件、會議記錄及進度追蹤
\item 參與用戶驗收測試及培訓安排
\item 協調跨部門溝通，確保項目里程碑達成
\end{itemize}
\textbf{關鍵字}：CRM、用戶驗收測試、項目協調、培訓}
\section{語言}

\cvline{粵語}{母語或雙語水平 \hfill \textbf{關鍵字}：母語}
\cvline{英語}{完全專業水平 \hfill \textbf{關鍵字}：IELTS 7.5、HKDSE English 5*}
\cvline{普通話}{完全專業水平 \hfill \textbf{關鍵字}：普通話水平測試二級甲等}
\cvline{日語}{初級水平 \hfill \textbf{關鍵字}：JLPT N4}
\section{專業技能}

\cvline{項目管理}{專家 \hfill \textbf{關鍵字}：敏捷開發、Scrum、瀑布模型、風險管理、持份者管理、預算控制}
\cvline{數碼轉型}{高級 \hfill \textbf{關鍵字}：流程自動化、雲端遷移、數據分析、變更管理}
\cvline{工具}{高級 \hfill \textbf{關鍵字}：Jira、Confluence、MS Project、Trello、Asana、Power BI}
\section{榮譽}

\cventry{2023年12月}
        {香港電訊數碼有限公司}
        {年度傑出項目經理}
        {}
        {}
        {表揚在數碼轉型項目中的卓越領導及交付表現，成功為公司節省 15\% 營運成本。}

\cventry{2020年11月}
        {新鴻基金融集團}
        {最佳創新項目獎}
        {}
        {}
        {領導的流動支付項目獲得集團年度創新獎，用戶採用率超過 80\%。}
\section{證書}

\cventry{2019年6月}
        {Project Management Institute}
        {專業項目管理師 (PMP)}
        {\url{https://www.pmi.org/certifications/project-management-pmp}}
        {}
        {}

\cventry{2021年9月}
        {Project Management Institute}
        {敏捷認證從業者 (PMI-ACP)}
        {\url{https://www.pmi.org/certifications/agile-acp}}
        {}
        {}

\cventry{2022年3月}
        {Amazon Web Services}
        {AWS Certified Cloud Practitioner}
        {\url{https://aws.amazon.com/certification/certified-cloud-practitioner/}}
        {}
        {}
\section{出版刊物}

\cventry{2023年5月}
        {香港企業數碼轉型的關鍵成功因素}
        {香港資訊科技期刊}
        {\url{https://www.hkitjournal.com/digital-transformation-hk}}
        {}
        {\begin{itemize}
\item 探討香港企業在數碼轉型過程中常見的挑戰與機遇
\item 提出項目管理在轉型中的關鍵角色及最佳實踐
\item 分析敏捷方法與傳統項目管理的混合應用
\end{itemize}}
\section{引薦}

\cventry{\emaillink[mankit.lee@ust.hk]{mankit.lee@ust.hk}}
        {香港科技大學資訊系統學系教授}
        {李文傑博士}
        {+852 9876 5432}
        {}
        {陳嘉敏在任職期間展現卓越的項目領導能力，能有效協調跨部門團隊，並在多個關鍵項目中達成超預期成果。}
\section{項目}

\cventry{2022年1月–2022年12月}
        {為大型電訊商建立整合式客戶服務平台}
        {智慧客戶服務平台}
        {\url{https://www.chan-ka-man.com/projects/smart-cs-platform}}
        {}
        {\begin{itemize}
\item 整合多渠道客服系統，包括電話、電郵、社交媒體及即時聊天
\item 導入 AI 聊天機械人，處理 60\% 常見查詢
\item 項目提前兩週上線，客戶滿意度提升 25\%
\end{itemize}
\textbf{關鍵字}：客戶服務、AI、系統整合}

\cventry{2019年3月–2019年11月}
        {為金融集團管理層開發實時業務數據儀表板}
        {金融數據儀表板}
        {\url{https://www.chan-ka-man.com/projects/finance-dashboard}}
        {}
        {\begin{itemize}
\item 整合多個內部系統數據，提供即時關鍵績效指標
\item 與業務分析師合作定義指標及可視化需求
\item 採用敏捷疊代方式交付，每兩週發布新功能
\end{itemize}
\textbf{關鍵字}：數據可視化、Power BI、KPI}
\section{興趣愛好}

\cvline{運動}{跑步、游泳、遠足}
\cvline{音樂}{鋼琴、爵士樂}
\cvline{志願服務}{青年 mentorship、社區服務}
\section{志願服務}

\cventry{2020年1月–2023年12月}
        {項目顧問}
        {香港青年協會}
        {\url{https://www.hkfyg.org.hk}}
        {}
        {\begin{itemize}
\item 為青年創業項目提供項目管理及商業策劃指導
\item 協助舉辦創新創業比賽及工作坊
\item 分享數碼轉型及職涯發展經驗
\end{itemize}}
    
\end{document}